\section{Conclusion}

In this work, we answer several scaling questions on training multilingual models, validating our findings on scales up to 8B model parameters. First, we introduce a novel scaling law formulation called Adaptive Transfer Scaling Law (ATLAS) that significantly improves scaling law fit by more than 0.3 $R^{2}$ for multilingual mixtures. To aid our analysis, we also derive transfer scores for 38 langauages. We find that English is often a strong positive transfer language, alongside languages with shared script and language families, though we note that this transfer is often asymmetric. However, there is still a compute-efficiency tax when we train multilingual models. We quantify this "curse of multilinguality" and provide a mathematical trade-off between model capacity and language coverage: we find that the performance penalty from adding languages is mitigated more effectively by scaling model size ($N$) rather than training tokens ($D$). Finally, we quantitatively answer a critical practical question for practitioners: determining the compute crossover point between pretraining and finetuning. We find that depending on the language, for a budget under 144-283B tokens, finetuning a general-purpose multilingual checkpoint is more efficient. Collectively, these contributions provide valuable guidelines for practitioners determining data mixtures, language coverage and computational requirements for creating multilingual foundation models.
